\section{Buist v. Anthropic — 「協調減速」独禁法訴訟の報道}
有料ユーザー4名が ChatGPT、Claude、Grok、Gemini の開発元を相手に、AI開発の協調減速が Sherman Act §1 に違反すると主張するクラスアクションを起こしたとの報道が拡散した。提訴日は9月18日とされ観測窓外だが、窓内で議論が広がったため収録する。
\dailyxsource{Bull Theory (@BullTheoryio)}{https://x.com/BullTheoryio/status/2101700126138061191}

\section{Claude Opus 5.5 とされるステルス試験リーク}
リーク投稿は、Anthropic が \texttt{claude-opus-5-5} を \texttt{claude-wafer-eap} として試験中で火曜公開予定だと主張。価格やthinking、tool useの仕様も共有されたが、観測窓内に公式確認はない。
\dailyxsource{lyra (@lyraxana)}{https://x.com/lyraxana/status/2101668658711798083}
\dailyxnote{未検証のリーク情報であり、仕様・公開時期を確定情報として扱わない。}

\section{GPT-6 系列の実務評価と次モデル観測}
Mikhail Parakhin は GPT-6 Pro を Math、ML、brainstorming で高く評価しつつ、Astra Max が勝るケースもあると報告した。OpenAI従業員の発言を GPT-6 Sol / Luna への暗示と読む解釈も流通したが、新モデルの公式発表ではない。
\dailyxsource{Mikhail Parakhin (@MParakhin)}{https://x.com/MParakhin/status/2101745841421861265}

\section{GPT-6 Astra の Minecraft エージェント報道}
IGN は、GPT-6 Astra の Minecraft エージェントが Creeper に進捗を壊された後、長時間ジャガイモを耕作し、緑色の対象を警戒するような挙動を示したと報道。長時間エージェントの失敗モードとして拡散した。
\dailyxsource{IGN (@IGN)}{https://x.com/IGN/status/2101748501763731672}

\section{Hugging Face買収後の中立性をめぐる議論}
SemiAnalysis は、NVIDIAによるHugging Face買収後の中立性に懸念を示し、一部ワークロードのModelScope移行検討に言及した。HF CEOのClemは、公開・中立・シリコン非依存へのコミットを示して反論した。
\dailyxsource{SemiAnalysis (@SemiAnalysis\_)}{https://x.com/SemiAnalysis_/status/2101703331286507887}

\section{Jensen Huang、既存法の適用を優先する規制観を提示}
NVIDIA公式は Jensen Huang の CBS Sunday 出演を紹介。「速く進むことと安全に作ることは対立しない」とし、関連クリップではAIラボと既存法の関係についての発言が共有され、規制論として議論を集めた。
\dailyxsource{NVIDIA (@nvidia)}{https://x.com/nvidia/status/2101683275634712630}

\section{dQwen3.5 — Hybrid-Attention Diffusion Language Models}
研究紹介投稿では、ハイブリッド注意とRNN系モデルを拡散言語モデルへ適用する dQwen3.5 を紹介。フル注意の対照モデルより同程度の損失へ少ないトークンで到達し、並列デコードにも強いとの主張が示された。
\dailyxsource{arXiv Bangers (@arXivBangers)}{https://x.com/arXivBangers/status/2101451318321877219}

\section{GitHub、本番投入前のLLM評価プラクティスを紹介}
GitHub公式は、クリーンなベンチマークで良好でも実運用では失敗し得るとして、プロトタイプから本番へ移す際のLLM評価実務を紹介。実際のユースケースを評価へ持ち込む必要性を強調する話題として観測された。
\dailyxsource{GitHub (@github)}{https://x.com/github/status/2101454543837929615}

\section{xAI / Grok、マーケター向けグループチャットを展開}
Josh Kim は、Grokの @bot を使うマーケター向けグループチャットを立ち上げ、Vercel、ClickUp、MongoDB、Framer、Perplexity、Zendesk、Bland の成長担当者が参加すると説明した。
\dailyxsource{Josh Kim (@joshkim)}{https://x.com/joshkim/status/2101763382085091817}

\section{Yann LeCun、AMI LabsとSSI、自動運転レベルを論評}
Yann LeCun は AMI Labs と SSI の調達規模や時期を比較し、現時点では双方とも製品リリースがないとの趣旨を投稿。別投稿では Waymo をLevel 4、Tesla FSDを「せいぜいLevel 3」と整理した。
\dailyxsource{Yann LeCun (@ylecun)}{https://x.com/ylecun/status/2101783007204114746}

\section{イランの学校攻撃とAI依存をめぐるBloomberg報道の再拡散}
窓内投稿は、欠陥インテリジェンス、古い画像、AIへの過度な依存がイランの学校攻撃に寄与したとする Bloomberg の調査報道を引用した。原報道は9月18日で観測窓外だが、窓内の安全保障議論として再拡散した。
\dailyxsource{I'MYOURHUCKLEBERRY (@PLSgetserious)}{https://x.com/PLSgetserious/status/2101791871651414454}

\section{エージェント経済、プライバシー、電力需要をめぐる議論}
Alex Finn は、エージェント競争の勝者がデータ・取引・広告を握るとの見方を提示。別系統ではWIREDのデータセンター需要報道や、AstraとFableをゲーム上で比較する配信も話題となった。
\dailyxsource{Alex Finn (@AlexFinn)}{https://x.com/AlexFinn/status/2101761022747140162}

\section{DeepSeek-v4.1-Flash のローカル推論トークン量報告}
投稿者はローカル環境で24時間に13億トークンを生成したとし、DeepSeek-v4.1-Flash は ML タスクで GLM-5.3-Flash より多く生成すると報告した。公式リリースとの対応や測定条件は未確認の自己報告である。
\dailyxsource{0xSero (@0xSero)}{https://x.com/0xSero/status/2101783078209233376}

\section{「来月のモデル一覧」コミュニティ予想}
コミュニティ投稿は Claude Opus 5.5、Gemini 4 Pro、GPT-6 Sol、Grok 4.8 / 4.7、Kimi K3.1、GLM 5.5 などをリリース候補として列挙した。公式カレンダーではなく、APIルーティング観測などを含む予想の集合である。
\dailyxsource{0xMarioNawfal (@RoundtableSpace)}{https://x.com/RoundtableSpace/status/2101736900801519713}
